\section*{Figure captions}
% generated by make_submission.py; do not edit by hand
\noindent\textbf{Figure 1.} How the dataset was built. $n$ counts scans, or corrections, at each step.

\noindent\textbf{Figure 2.} The two anchors, from the released labels. Red, lowest rib-bearing vertebra and its rib; blue, S1; yellow, the rib-free vertebrae between. (b) and (c) both count four. In (b) a long rib sits on L1; in (c) the twelfth ribs are stumps and the segment below the fourth lumbar body is fused to the sacrum (Castellvi IIIb) and carries the S1 identifier.

\noindent\textbf{Figure 3.} Instrumentation in the release, one exemplar each, drawn through bone. (a) Hip arthroplasty, 64, $n=8$. (b) Femoral-neck screws, 66, $n=1$. (c) Sacroiliac screws, 65, $n=1$. (d) Interbody cages, 61, $n=1$. Bar, 5\,cm.

\noindent\textbf{Figure 4.} Count-free measures. (a) Vertebrae between the lowest rib and the sacrum. (b) Lowest rib length over the rib above. (c) Lowest lumbar transverse span against its gap to the ala; transitional labels marked.

\noindent\textbf{Figure 5.} Spinopelvic measures against standing reference bands (mean $\pm$ SD)~\cite{vialle2005}.

\noindent\textbf{Figure 6.} Morphometry by level, as median, interquartile range and 5th--95th percentile, with $n$ at right. (a) Body height. (b) Endplate width. (c) Canal. (d) Pedicle width. (e) Disc height. (f) Trabecular attenuation.

\noindent\textbf{Figure 7.} Records carrying each vertebral level, cranial to caudal. Counts per identifier in Table~S1.

